\section{Líneas de trabajo futuro}
\label{sec:lineas_futuras}

El desarrollo de Voctomix 2.0 ha puesto de manifiesto varias áreas de mejora e
investigación que quedan fuera del alcance de este Trabajo de Fin de Grado pero que
constituyen líneas de trabajo futuro de interés:

\begin{itemize}
  \item \textbf{Prueba con cámaras reales:} la validación definitiva del sistema
        consiste en conectar cámaras físicas a los puertos de entrada de voctocore
        y realizar una retransmisión en directo completa. Este caso de uso está
        previsto en el contexto del grupo de investigación CyberNemo de la UPM,
        donde Voctomix 2.0 se utilizará para retransmitir seminarios y presentaciones
        con cámaras reales conectadas mediante tarjetas de captura o directamente
        mediante NDI. Esta prueba permitirá validar el sistema en condiciones reales
        y detectar posibles problemas de sincronización, latencia o estabilidad que
        no se manifiestan con fuentes sintéticas.

  \item \textbf{Interfaz web de control:} sustituir voctogui por una aplicación web
        que consuma la API REST del servicio de telemetría eliminaría la dependencia
        de GTK y permitiría controlar el sistema desde cualquier dispositivo con
        navegador, incluidas tabletas y teléfonos móviles.

  \item \textbf{Soporte NDI nativo:} incorporar el protocolo NDI como fuente de
        entrada alternativa a TCP/MKV facilitaría la integración con cámaras de red
        modernas y con herramientas de producción compatibles con este protocolo.
\end{itemize}
